\documentclass[11pt, letterpaper]{article}
\usepackage{amsmath, amssymb}
\usepackage{hyperref}
\usepackage{booktabs}
\usepackage{enumitem}
\usepackage{xcolor}
\usepackage{geometry}
\usepackage{fancyhdr}
\usepackage{titlesec}
\usepackage{tcolorbox}
\usepackage{graphicx}
\usepackage{parskip}
\usepackage{caption}
\geometry{margin=1in}

\definecolor{courseblue}{HTML}{002452}
\definecolor{coursegold}{HTML}{FABD0F}
\definecolor{coursered}{HTML}{B90E31}
\definecolor{lightblue}{HTML}{E8EEF7}
\definecolor{lightgold}{HTML}{FEF8E3}

\hypersetup{colorlinks=true,linkcolor=courseblue,urlcolor=coursered}
\tcbuselibrary{skins}
\newtcolorbox{guidancebox}{enhanced,colback=lightblue,colframe=courseblue,
  sharp corners,leftrule=4pt,left=6pt,right=6pt,top=4pt,bottom=4pt}
\newtcolorbox{notebox}{enhanced,colback=lightgold,colframe=coursegold,
  sharp corners,leftrule=4pt,left=6pt,right=6pt,top=4pt,bottom=4pt}

\pagestyle{fancy}\fancyhf{}
\fancyhead[L]{\small\textcolor{courseblue}{\textbf{CISC 867 --- Deep Learning | Fall 2026}}}
\fancyhead[R]{\small\textcolor{courseblue}{Midterm Report}}
\fancyfoot[C]{\small\thepage}
\renewcommand{\headrulewidth}{0.4pt}
\titleformat{\section}{\large\bfseries\color{courseblue}}{}{0em}{}[\titlerule]
\titleformat{\subsection}{\normalsize\bfseries\color{courseblue}}{}{0em}{}

\begin{document}

{\centering
  {\Large\textbf{CISC 867: Deep Learning --- Midterm Report}}\\[0.5em]
  {\large\textit{[PROJECT ID AND TITLE]}}\\[0.4em]
  \textbf{Group Members:}\\
  {[Student 1 Name]} (Project Lead) \quad
  {[Student 2 Name]} (Research Lead) \quad
  {[Student 3 Name]} (Engineering Lead)\\[0.3em]
  \textbf{GitHub:} \url{[link]} \quad \textbf{Date:} \today\\[0.4em]
  \rule{\linewidth}{0.4pt}
}

\section{Abstract}

\begin{guidancebox}
\textbf{Guidance.} 150--200 word summary of: (1) the problem, (2) what you have
implemented so far, (3) your key preliminary result, and (4) your planned next steps.
\end{guidancebox}

[WRITE YOUR ABSTRACT HERE.]

\section{Introduction and Problem Statement}

\begin{guidancebox}
\textbf{Guidance.} Brief recap of the problem and why it matters (1/2 page max).
Focus on what has changed since the proposal---have you refined the scope or identified
a clearer novel angle? Include the \textbf{research question} you are now testing.
\end{guidancebox}

[CONCISE PROBLEM DESCRIPTION AND RESEARCH QUESTION. NOTE ANY SCOPE CHANGES.]

\section{Related Work (Updated)}

\begin{guidancebox}
\textbf{Guidance.} An extended version of your proposal survey. Add papers you
discovered since the proposal. Include a comparative table or taxonomy where possible.
Make explicit connections to your own approach. Add at least 2 new papers beyond
the proposal. Target: $\approx$400--600 words.
\end{guidancebox}

[UPDATED LITERATURE REVIEW.]

\begin{center}
\begin{tabular}{llp{4cm}p{4cm}}
\toprule
\textbf{Method} & \textbf{Year} & \textbf{Key idea} & \textbf{Relation to our work} \\
\midrule
{[Method 1]} & [Year] & [Key idea] & [Relation] \\
{[Method 2]} & [Year] & [Key idea] & [Relation] \\
{[Method 3] }& [Year] & [Key idea] & [Relation] \\
\bottomrule
\end{tabular}
\end{center}

\section{Methodology}

\begin{guidancebox}
\textbf{Guidance.} Describe what you have implemented in detail. Include:
(1) architecture / model description with any modifications,
(2) training procedure, hyperparameters, and optimiser,
(3) data preprocessing pipeline,
(4) evaluation protocol (metric definitions and dataset splits).
Use equations where appropriate.
\end{guidancebox}

\subsection{Model Architecture}

[DESCRIBE YOUR MODEL. If you are using a pre-trained backbone, state clearly
which weights you use and what you add/modify on top.]

\subsection{Training Setup}

\begin{tabular}{ll}
\toprule
\textbf{Hyperparameter} & \textbf{Value} \\
\midrule
Optimiser & [e.g., AdamW] \\
Learning rate & [e.g., $1\times10^{-4}$] \\
Batch size & [e.g., 32] \\
Epochs / iterations & [e.g., 50 epochs] \\
Loss function & [e.g., Cross-entropy + Dice] \\
Data augmentation & [e.g., Random flip, colour jitter] \\
Random seeds used & [e.g., 42, 123, 999] \\
\bottomrule
\end{tabular}

\subsection{Dataset and Preprocessing}

[DESCRIBE THE DATASET SPLIT (train/val/test sizes), ANY PREPROCESSING STEPS
(resizing, normalisation, augmentation), AND HOW YOU HANDLE CLASS IMBALANCE
IF APPLICABLE.]

\section{Preliminary Results}

\begin{guidancebox}
\textbf{Guidance.} Present quantitative results for your baseline (and, if
available, early results for your proposed method). Report mean $\pm$ std over
at least 3 seeds. Compare to at least one published baseline. Include at least
one qualitative example (figure or table). Discuss what the numbers tell you.
\end{guidancebox}

\subsection{Baseline Reproduction}

\begin{notebox}
\textbf{Target:} Match the published number within 1--2\% before proceeding
to your proposed modification.
\end{notebox}

\begin{center}
\begin{tabular}{lcc}
\toprule
\textbf{Model} & \textbf{Metric 1 (Mean $\pm$ Std)} & \textbf{Metric 2 (Mean $\pm$ Std)} \\
\midrule
Published baseline & 3.17 &  9.46$\pm$ 0.11 \\
Our reproduction & 5.71 & [Your value] $\pm$ [Std] \\
Our proposed (early) & N/A & N/A \\
\bottomrule
\end{tabular}
\captionof{table}{Preliminary results on MNIST, 10,000 images. All results are mean $\pm$ std over 3 seeds.}

INSERT TABLE HERE

\end{center}

\subsection{Qualitative Analysis}

[INSERT A REPRESENTATIVE FIGURE: e.g., a before/after image pair, a learning
curve, a confusion matrix, or a prediction example. Caption should fully
describe what is being shown.]

\begin{figure}[h]
  \centering
  % \includegraphics[width=0.8\linewidth]{figures/qualitative.png}
  \caption{Comparison between a DDPM generated set of images, and the MNIST images the model was trained on. On a 28x28 grid, the model's output looks indistinguishable from the training dataset's images that were drawn by hand.}
  \label{fig:qual}
\end{figure}

\subsection{Discussion}

[INTERPRET YOUR RESULTS. What works? What does not? What surprised you?
Where are the biggest failure cases? What do you conclude so far?]

\section{Challenges and Adaptations}

\begin{guidancebox}
\textbf{Guidance.} Describe any obstacles you encountered and how you
addressed or plan to address them. Did you deviate from your proposal timeline?
If so, explain why and provide an updated plan.
\end{guidancebox}

There were quite a few technical obstacles with the DDPM and DDIM versions of the project, mainly centered around compute power. Google Colab's free version sometimes did not allow us to access the GPU's and TPU's needed to train our model, which caused issues in trying to run the model multiple times to capture results. When we had access to the necessary compute, the datasets and model mostly acted as intended and allowed us to experiment and work on the model freely.
Developing the DDPM was difficult, especially conceptually understanding the build of a UNet and how a trained UNet would be able to generate an image, but working more with PyTorch and understanding the literature from Ho et. Al (INSERT REFERENCE HERE) helped tremendously. After the initial concept of the DDPM was created and working as intended, building off it for the different variations and the DDIM was much easier.
Going forward, we plan to run all tests for the finished version of each model one after the other, to avoid being locked out of GPU/TPU use by Google Colab. We also intend to dive deeply into the Classifier Free Guidance and Consistency model literature before we start building, so we have a greater understanding of the architecture before we start building.

\subsection{Updated Week 7--10 Plan}

\begin{tabular}{lp{8cm}l}
\toprule
\textbf{Week} & \textbf{Tasks} & \textbf{Responsible} \\
\midrule
7 & Build and test Classifier Free Guidance model & Spencer \\
8 & Build and test Consistency model, start final report & Oliver \\
9 & Finish final report and practice presentation & All \\
10 & Submit final report and present \textbf{(Final due)} & All \\
\bottomrule
\end{tabular}

\section{Preliminary Novel Contribution (Optional Bonus)}

\begin{guidancebox}
\textbf{Guidance.} If you have a preliminary finding that goes beyond
reproducing the baseline---an unexpected result, an improved ablation,
an insight about the method---describe it here for the +2 bonus.
Provide quantitative evidence.
\end{guidancebox}

[DESCRIBE YOUR PRELIMINARY NOVEL FINDING, WITH EVIDENCE. OR LEAVE EMPTY.]

\section{Contribution Declaration}

\begin{tabular}{lcc}
\toprule
\textbf{Member} & \textbf{Contributions (this period)} & \textbf{\%} \\
\midrule
{[Student 1]} & [e.g., Lit review, updated proposal section] & \% \\
{[Student 2]} & [e.g., Model implementation, training runs] & \% \\
{[Student 3] }& [e.g., Data pipeline, evaluation code] & \% \\
\bottomrule
\end{tabular}

\section{References}

\begin{enumerate}[label={[\arabic*]},leftmargin=*]
  \item [Author(s), Year. ``Title.'' \textit{Venue}. arXiv:XXXX.XXXXX]
  \item [\ldots]
\end{enumerate}

\end{document}